% !TEX root = ../../ctfp-print.tex

\lettrine[lhang=0.17]{现}{在你已经知道}什么是函子（functor），并且看过一些例子之后，我们来看看如何从较小的函子构建更大的函子。特别是值得研究哪些类型构造器（对应范畴间对象的映射）可以扩展为包含态射间映射的函子。

\section{双函子}

由于函子是范畴范畴 $\Cat$ 中的态射，关于态射（尤其是函数）的许多直观认识同样适用于函子。例如，就像存在接受两个参数的函数一样，我们也存在接受两个参数的\newterm{双函子（bifunctor）}。在对象层面，双函子将来自范畴 $\cat{C}$ 和范畴 $\cat{D}$ 的每对对象分别映射到范畴 $\cat{E}$ 中的对象。注意这实际上就是说它是从范畴的\newterm{笛卡尔积（Cartesian product）} $\cat{C}\times{}\cat{D}$ 到 $\cat{E}$ 的映射。

\begin{figure}[H]
  \centering\includegraphics[width=0.3\textwidth]{images/bifunctor.jpg}
\end{figure}

\noindent
这很直观明了。但函子性要求双函子必须同时映射态射。这次它需要将一对态射（一个来自 $\cat{C}$，一个来自 $\cat{D}$）映射为 $\cat{E}$ 中的态射。

同样地，这对态射本质上就是乘积范畴 $\cat{C}\times{}\cat{D}$ 中的单个态射。我们定义范畴笛卡尔积中的态射为从一对对象到另一对对象的态射对。这些态射对显然可组合：
\[(f, g) \circ (f', g') = (f \circ f', g \circ g')\]
这种组合满足结合律且存在单位元——即恒等态射对 $(\id, \id)$。因此范畴的笛卡尔积确实构成范畴。

理解双函子的另一种方式是将其视为每个参数分别独立的函子。因此与其将函子律（结合性和单位元保持）从函子推广到双函子，只需分别验证每个参数成立即可。然而一般而言，单独的函子性不足以证明联合函子性。那些联合函子性不成立的范畴被称为\newterm{预单子（premonoidal）}。

让我们在Haskell中定义双函子。此时三个范畴完全相同：都是Haskell类型范畴。双函子是接受两个类型参数的类型构造器。以下是直接取自库文件 \code{Control.Bifunctor} 的 \code{Bifunctor} 类型类定义：

\src{snippet01}

\begin{figure}[H]
  \centering\includegraphics[width=0.3\textwidth]{images/bimap.jpg}
  \caption{bimap}
\end{figure}

类型变量 \code{f} 表示双函子。可以看到在所有类型签名中它总是应用于两个类型参数。第一个类型签名定义了 \code{bimap}：同时映射两个函数。结果是一个提升后的函数 \code{(f a b -> f c d)}，作用于双函子类型构造器生成的类型。该类型类提供了通过 \code{first} 和 \code{second} 实现 \code{bimap} 的默认版本。（如前所述，这种方法并不总是有效，因为两次映射可能不满足交换律，即 \code{first g . second h} 可能不等于 \code{second h . first g}。）

\begin{figure}[H]
  \centering
  \begin{minipage}{0.45\textwidth}
    \centering
    \includegraphics[width=0.65\textwidth]{images/first.jpg} % first figure itself
    \caption{first}
  \end{minipage}\hfill
  \begin{minipage}{0.45\textwidth}
    \centering
    \includegraphics[width=0.6\textwidth]{images/second.jpg} % second figure itself
    \caption{second}
  \end{minipage}
\end{figure}

\noindent
另外两个类型签名 \code{first} 和 \code{second} 分别见证了 \code{f} 在第一个和第二个参数上的函子性。

类型类定义提供了基于 \code{bimap} 的这两个方法的默认实现。

当声明 \code{Bifunctor} 实例时，可以选择两种方式：要么实现 \code{bimap} 并接受 \code{first} 和 \code{second} 的默认实现；或者实现 \code{first} 和 \code{second} 并接受 \code{bimap} 的默认实现（当然也可以全部三个都实现，但需要自行保证它们之间的关系）。

\section{乘积与余积双函子}

一个重要的双函子例子是范畴论中的乘积——由\hyperref[products-and-coproducts]{泛构造}定义的两个对象的乘积。若范畴中任意两个对象都存在乘积，则从对象到其乘积的映射具有双函子性质。这在一般范畴论中成立，在Haskell语言中尤其如此。以下是针对序对构造器的\code{Bifunctor}实例——最简单的乘积类型：

\src{snippet02}[b]
这里的实现别无选择：\code{bimap}函数只能分别将第一个函数作用于序对的第一个分量，第二个函数作用于第二个分量。给定类型签名后，代码几乎是自动生成的：

\src{snippet03}
该双函子的作用是生成类型对，例如：

\begin{snip}{haskell}
(,) a b = (a, b)
\end{snip}
通过对偶原理，若范畴中每个对象对都有余积（coproduct），则余积也构成双函子。在Haskell中，\code{Either}类型构造器作为\code{Bifunctor}的实例即体现了这一点：

\src{snippet04}[b]
这段代码同样可以自动推导。

回想我们讨论幺半范畴时的情形：幺半范畴定义了作用于对象的二元运算符和单位对象。我曾提到$\Set$范畴关于笛卡尔积构成幺半范畴，其单位元是单元素集合；它也关于不交并构成幺半范畴，此时单位元是空集。当时未提及的一个重要条件是：该二元运算符必须是双函子。这是幺半范畴定义的关键要求——我们要求幺半积与范畴结构（由态射定义）相容。现在我们已更接近完整的幺半范畴定义（仍需学习自然性概念才能最终完成定义）。

\section{函子代数数据类型}

我们已经见过多个参数化数据类型的例子，它们最终都成为了函子---我们可以为它们定义\code{fmap}。复杂的数据类型由更简单的数据类型构建而成。特别是代数数据类型（\acronym{ADT}）通过求和与乘积构造而成。我们刚刚看到求和与乘积具有函子性质。我们也知道函子可以复合。因此若能证明\acronym{ADT}的基本构件具有函子性质，则参数化\acronym{ADT}必然也具有函子性质。

那么参数化代数数据类型的基本构件是什么？首先是那些不依赖函子类型参数的项，例如\code{Maybe}中的\code{Nothing}，或\code{List}中的\code{Nil}。它们等价于\code{Const}函子。记住，\code{Const}函子会忽略其类型参数（严格来说是第二个类型参数，这是与我们相关的参数，第一个参数保持不变）。

其次是直接封装类型参数本身的元素，例如\code{Maybe}中的\code{Just}。它们等价于恒等函子。我之前提到过恒等函子作为范畴\emph{Cat}中的恒等态射，但尚未给出其在Haskell中的定义。现在展示如下：

\src{snippet05}

\src{snippet06}
你可以将\code{Identity}视为最简单的容器，它始终存储一个类型为\code{a}的（不可变）值。

代数数据结构中的其他所有内容都可以通过这两个基本构件用乘积和求和来构造。

有了这个新知识后，让我们重新审视\code{Maybe}类型构造器：

\src{snippet07}
它是两个类型的求和，我们现在知道求和具有函子性质。第一部分\code{Nothing}可以表示为作用于\code{a}的\code{Const ()}（\code{Const}的第一个类型参数设为单位类型）。第二部分正是恒等函子的不同名称。我们本可以将\code{Maybe}（直到同构）定义为：

\src{snippet08}
因此\code{Maybe}是双函子\code{Either}与两个函子\code{Const ()}和\code{Identity}的复合。（\code{Const}本质上是双函子，但这里我们总是部分应用它。）

我们已经知道函子的复合仍是函子---我们可以轻易验证双函子的情况同样成立。我们需要明确的是双函子与两个函子复合在态射上的作用方式。给定两个态射，我们只需分别用一个函子提升其中一个态射，另一个函子提升另一个态射。然后用双函子提升这对被提升的态射。

我们可以在Haskell中表达这种复合。定义一个数据类型，它以双函子\code{bf}（这是一个接受两个类型作为参数的类型构造器）为参数，两个函子\code{fu}和\code{gu}（每个都是接受一个类型变量的类型构造器），以及两个常规类型\code{a}和\code{b}。我们将\code{fu}应用于\code{a}，\code{gu}应用于\code{b}，然后将结果的两个类型应用于\code{bf}：

\src{snippet09}
这是对象（或类型）上的复合。注意在Haskell中我们如何应用类型构造器到类型上，就像函数应用到参数一样。语法是相同的。

如果你感到困惑，尝试按顺序将\code{BiComp}应用于\code{Either}、\code{Const ()}、\code{Identity}、\code{a}和\code{b}。你将得到我们简化版的\code{Maybe b}（\code{a}被忽略）。

新的数据类型\code{BiComp}是\code{a}和\code{b}上的双函子，但这仅当\code{bf}本身是\code{Bifunctor}而\code{fu}和\code{gu}是\code{Functor}时成立。编译器必须知道\code{bf}存在\code{bimap}定义，而\code{fu}和\code{gu}存在\code{fmap}定义。在Haskell中，这通过实例声明中的前置条件表达：一组类约束后接双箭头：

\src{snippet10}[b]
\code{BiComp}的\code{bimap}实现基于\code{bf}的\code{bimap}以及\code{fu}和\code{gu}的两个\code{fmap}。编译器会自动推导所有类型，并在使用\code{BiComp}时选择正确的重载函数。

定义中\code{bimap}的\code{x}具有以下类型：

\src{snippet11}
这相当复杂。外层\code{bimap}穿透外层的\code{bf}层，两个\code{fmap}分别深入\code{fu}和\code{gu}内部。若\code{f1}和\code{f2}的类型为：

\src{snippet12}
则最终结果类型为\code{bf (fu a') (gu b')}：

\src{snippet13}[b]
如果你喜欢拼图游戏，这类类型操作可能会带给你数小时的乐趣。

因此事实证明，我们无需单独证明\code{Maybe}是函子---这个结论自然地从它作为两个函子基本构件的求和构造中得出。

敏锐的读者可能会问：如果代数数据类型的\code{Functor}实例推导如此机械，能否自动化并由编译器完成？确实可以，而且实际上就是这样实现的。你需要在源文件顶部添加以下行启用特定Haskell扩展：

\begin{snip}{haskell}
{-# LANGUAGE DeriveFunctor #-}
\end{snip}
然后向你的数据结构添加\code{deriving Functor}：

\begin{snip}{haskell}
data Maybe a = Nothing | Just a deriving Functor
\end{snip}
相应的\code{fmap}将会自动为你实现。

代数数据结构的规律性使得不仅能够推导\code{Functor}实例，还能推导包括前文提到的\code{Eq}类型类在内的多个其他类型类实例。你也可以教导编译器推导你自己类型类的实例，但这稍微高级些。核心思想相同：你为基本构件、求和与乘积提供行为定义，让编译器自行处理其余部分。

\section{C++中的函子}

如果你是一位C++程序员，显然在实现函子(functor)方面只能依靠自己。然而，你应该能够识别C++中某些类型的代数数据结构(algebraic data structure)。当这样的数据结构被实现为泛型模板时，你应该能快速为其编写\code{fmap}实现。

让我们观察一个树(tree)数据结构，在Haskell中我们会将其定义为递归的和类型(sum type)：

\src{snippet14}
如前所述，C++中实现和类型的一种方式是通过类继承体系。在面向对象语言中，很自然会想到将\code{fmap}作为基类\code{Functor}的虚函数，并在所有子类中重写它。不幸的是这种方法不可行，因为\code{fmap}是一个模板函数，不仅参数化了操作对象的类型（即\code{this}指针），还参数化了应用函数的返回类型。C++不允许虚函数模板化。我们将把\code{fmap}实现为泛型自由函数，并用\code{dynamic\_cast}替代模式匹配。

基类必须至少定义一个虚函数才能支持动态转型，因此我们将析构函数设为虚函数（这本身也是良好实践）：

\begin{snip}{cpp}
template<class T>
struct Tree {
    virtual ~Tree() {}
};
\end{snip}
\code{Leaf}本质上就是伪装的\code{Identity}函子：

\begin{snip}{cpp}
template<class T>
struct Leaf : public Tree<T> {
    T _label;
    Leaf(T l) : _label(l) {}
};
\end{snip}
\code{Node}则是乘积类型(product type)：

\begin{snip}{cpp}
template<class T>
struct Node : public Tree<T> {
    Tree<T> * _left;
    Tree<T> * _right;
    Node(Tree<T> * l, Tree<T> * r) : _left(l), _right(r) {}
};
\end{snip}
实现\code{fmap}时我们利用了\code{Tree}类型的动态分发机制。\code{Leaf}情况应用\code{Identity}版本的\code{fmap}，而\code{Node}情况则视为双函子(bifunctor)与两个\code{Tree}函子的复合。作为C++程序员，你可能不太习惯用这些范畴论术语分析代码，但这有助于培养范畴论思维。

\begin{snip}{cpp}
template<class A, class B>
Tree<B> * fmap(std::function<B(A)> f, Tree<A> * t) {
    Leaf<A> * pl = dynamic_cast <Leaf<A>*>(t);
    if (pl)
        return new Leaf<B>(f (pl->_label));
    Node<A> * pn = dynamic_cast<Node<A>*>(t);
    if (pn)
        return new Node<B>( fmap<A>(f, pn->_left)
                          , fmap<A>(f, pn->_right));
    return nullptr;
}
\end{snip}
为简化起见，我忽略了内存和资源管理问题，但在生产代码中建议使用智能指针（根据策略选择unique或shared指针）。

对比Haskell的\code{fmap}实现：

\src{snippet15}
该实现还可以由编译器自动推导生成。

\section{Writer 函子}

我曾承诺会回到之前介绍的 \hyperref[kleisli-categories]{Kleisli
  范畴}。该范畴中的态射表现为返回 \code{Writer}
数据结构的「装饰」函数。

\src{snippet16}
我说过这种装饰与自函子（endofunctor）存在某种关联。事实上，
\code{Writer} 类型构造器在 \code{a} 上具有函子性（functorial）。我们甚至无需手动实现其 \code{fmap}，
因为它本质上是一个简单的乘积类型。

但 Kleisli 范畴与函子之间究竟有什么普遍联系呢？作为范畴，Kleisli 范畴定义了组合律和恒等态射。提醒您：
组合律由鱼操作符（fish operator）给出：

\src{snippet17}
而恒等态射由名为 \code{return} 的函数表示：

\src{snippet18}
有趣的是，当你长时间观察这两个函数的类型（我是说，足够长的时间），你会找到一种方法将它们组合起来，
生成一个具有正确类型签名的 \code{fmap} 函数。具体方式如下：

\src{snippet19}
这里，鱼操作符组合了两个函数：其中一个是我们熟悉的 \code{id}，另一个是将 \code{f} 作用于 lambda 参数后，
再用 \code{return} 处理结果的 lambda 表达式。最令人费解的部分可能是对 \code{id} 的使用。毕竟，
鱼操作符的输入难道不是应该接受一个接收「普通」类型并返回装饰类型的函数吗？其实不然。
没人规定 \code{a -> Writer b} 中的 \code{a} 必须是「普通」类型。由于它是类型变量，
它可以代表任何类型，特别是可以是装饰类型，比如 \code{Writer b}。

因此 \code{id} 将接收 \code{Writer a} 并输出 \code{Writer a}。鱼操作符会从中提取出
\code{a} 值，并将其作为参数 \code{x} 传入 lambda 表达式。在此处，\code{f} 将其转换为
\code{b}，然后 \code{return} 对其进行装饰，最终得到 \code{Writer b}。整体来看，
我们得到了一个接收 \code{Writer a} 并返回 \code{Writer b} 的函数——这正是 \code{fmap}
应当产生的结果。

值得注意的是，这个论证具有普适性：你可以将 \code{Writer} 替换为任意类型构造器。只要它支持鱼操作符和
\code{return}，就能定义相应的 \code{fmap}。因此 Kleisli 范畴中的装饰必然构成一个函子。（不过并非所有函子都能导出 Kleisli 范畴。）

你可能会疑惑：我们刚刚定义的 \code{fmap} 是否与编译器通过 \code{deriving Functor} 自动生成的版本一致？
有趣的是，答案是肯定的。这源于 Haskell 实现多态函数的方式，称为参数多态性（parametric polymorphism），
它会产生所谓的「免费定理」（theorems for free）。其中一个定理指出：若某个类型构造器存在保持恒等性的
\code{fmap} 实现，则该实现必然是唯一的。

\section{协变与逆变函子}

回顾了writer函子后，我们回到reader函子。它基于部分应用的函数箭头类型构造器：

\src{snippet20}
我们可以将其重写为类型同义词：

\src{snippet21}
其对应的\code{Functor}实例正如之前所见：

\src{snippet22}
但与pair类型构造器或\code{Either}类型构造器类似，函数类型构造器也接受两个类型参数。pair和\code{Either}在其两个参数上都是函子性的——它们是双函子(bifunctor)。那么函数类型构造器是否也是双函子呢？

让我们尝试使其在第一个参数上具有函子性。从一个类型同义词开始——这类似于\code{Reader}但参数顺序反转：

\src{snippet23}
这次我们固定返回类型\code{r}，改变参数类型\code{a}。能否设法匹配类型来实现\code{fmap}，其类型签名如下：

\src{snippet24}
仅凭两个分别接收\code{a}并返回\code{b}和\code{r}的函数，我们根本无法构造出接收\code{b}并返回\code{r}的函数！如果能以某种方式反转第一个函数，使其接收\code{b}并返回\code{a}就会不同。虽然我们无法反转任意函数，但可以转向对偶范畴(opposite category)。

简要回顾：对于每个范畴$\cat{C}$都存在对偶范畴$\cat{C}^\mathit{op}$。这个范畴与$\cat{C}$有相同对象，但所有箭头方向反转。

考虑一个从$\cat{C}^\mathit{op}$到其他范畴$\cat{D}$的函子：
$$F \Colon \cat{C}^\mathit{op} \to \cat{D}$$
该函子将$\cat{C}^\mathit{op}$中的态射$f^\mathit{op} \Colon a \to b$映射到$\cat{D}$中的态射$F f^\mathit{op} \Colon F a \to F b$。但$f^\mathit{op}$实际上对应于原范畴$\cat{C}$中的某个态射$f \Colon b \to a$。注意这里的反转。

现在，$F$是一个常规函子，但我们还可以基于$F$定义另一种映射——称之为$G$。这个映射不是函子，它从$\cat{C}$到$\cat{D}$。它以相同方式映射对象，但在映射态射时会反转方向。它取$\cat{C}$中的态射$f \Colon b \to a$，先将其映射为对偶态射$f^\mathit{op} \Colon a \to b$，再应用函子$F$得到$F f^\mathit{op} \Colon F a \to F b$。

考虑到$F a$与$G a$相同且$F b$与$G b$相同，整个过程可描述为：$G f \Colon (b \to a) \to (G a \to G b)$
这是一种"带有扭曲的函子"。这种以反转态射方向方式映射范畴的结构称为\emph{逆变函子}(contravariant functor)。注意逆变函子本质上就是从对偶范畴出发的常规函子。顺便提一下，我们迄今为止研究的常规函子被称为\emph{协变}(covariant)函子。

\begin{figure}[H]
  \centering
  \includegraphics[width=40mm]{images/contravariant.jpg}
\end{figure}

\noindent
以下是Haskell中定义逆变函子(严格来说是逆变\emph{自函子})的类型类：

\src{snippet25}
我们的类型构造器\code{Op}是其实例：

\src{snippet26}
注意函数\code{f}插入到\code{Op}内容(即函数\code{g})的\emph{前面}(右侧)。

若注意到\code{contramap}的定义本质上是反转参数顺序的函数组合操作符，我们可以用更简洁的方式定义\code{Op}的实例。标准库提供了一个专门反转参数顺序的函数\code{flip}：

\src{snippet27}
利用它我们可以得到：

\src{snippet28}

\section{普罗函子（Profunctors）}

我们已经看到，函数箭头运算符在其第一个参数中是逆变（contravariant）的，在第二个参数中是协变（covariant）的。这种结构有专门的名称吗？事实上，如果目标范畴是$\Set$，这类结构被称为\newterm{普罗函子（profunctor）}。由于逆变函子等价于从对偶范畴出发的协变函子，普罗函子被定义为如下形式的函子：
$$\cat{C}^\mathit{op} \times \cat{D} \to \Set$$
鉴于Haskell类型可近似视为集合，我们将名称\code{Profunctor}赋予具有两个类型参数的类型构造子\code{p}，它在第一个参数中逆变，在第二个参数中协变。以下是来自\code{Data.Profunctor}库的相关类型类定义：

\src{snippet29}[b]
这三个函数均提供默认实现。与\code{Bifunctor}类似，在声明\code{Profunctor}实例时，可以选择实现\code{dimap}并接受\code{lmap}和\code{rmap}的默认实现，或分别实现\code{lmap}和\code{rmap}并接受\code{dimap}的默认实现。

\begin{figure}[H]
  \centering
  \includegraphics[width=0.4\textwidth]{images/dimap.jpg}
  \caption{dimap}
\end{figure}

\noindent
现在我们可以断言函数箭头运算符是\code{Profunctor}的一个实例：

\src{snippet30}[b]
普罗函子在Haskell的Lens库中具有重要应用。当我们讨论端（ends）与余端（coends）时会再次遇到它们。

\section{同态函子}

上述例子反映了一个更普遍的命题：将对象对$a$和$b$映射到它们之间的态射集合（即同态集$\cat{C}(a, b)$）的映射是一个函子。这个函子从积范畴$\cat{C}^\mathit{op}\times{}\cat{C}$到集合范畴$\Set$。

我们定义它在态射上的作用。积范畴$\cat{C}^\mathit{op}\times{}\cat{C}$中的态射是一对来自$\cat{C}$的态射：
\begin{gather*}
  f \Colon a' \to a \\
  g \Colon b \to b'
\end{gather*}
这对态射的提升应当是从集合$\cat{C}(a, b)$到集合$\cat{C}(a', b')$的一个态射（即函数）。任取$\cat{C}(a, b)$中的元素$h$（它是从$a$到$b$的态射），将其映射为：
$$g \circ h \circ f$$
这显然是$\cat{C}(a', b')$中的一个元素。

如你所见，同态函子是缩影函子（profunctor）的一个特例。

\section{挑战}

\begin{enumerate}
  \tightlist
  \item
        证明以下数据类型：

        \begin{snip}{haskell}
data Pair a b = Pair a b
\end{snip}

        是一个双函子（bifunctor）。进阶任务：实现\code{Bifunctor}的所有三个方法，并通过等式推理证明这些定义在可应用时与默认实现兼容。
  \item
        展示标准定义的\code{Maybe}类型与此脱糖表示之间的同构性：

        \begin{snip}{haskell}
type Maybe' a = Either (Const () a) (Identity a)
\end{snip}

        提示：定义两种实现间的双向映射。进阶任务：使用等式推理证明它们互为逆函数。
  \item
        尝试另一种数据结构。我称其为\code{PreList}，因为它是\code{List}的雏形。它用类型参数\code{b}替代递归：

        \begin{snip}{haskell}
data PreList a b = Nil | Cons a b
\end{snip}

        通过递归地将\code{PreList}应用于自身即可恢复早先的\code{List}定义（当我们讨论不动点时会看到具体实现方式）。

        证明\code{PreList}是\code{Bifunctor}的一个实例。
  \item
        证明以下数据类型在\code{a}和\code{b}中定义了双函子：

        \begin{snip}{haskell}
data K2 c a b = K2 c

data Fst a b = Fst a

data Snd a b = Snd b
\end{snip}

        进阶任务：将你的解答与Conor McBride的论文《Clowns to the Left of me, Jokers to the Right》进行对照验证 \urlref{http://strictlypositive.org/CJ.pdf}{(点击访问原文)}。
  \item
        在非Haskell语言中定义双函子。为该语言中的泛型对实现\code{bimap}函数。
  \item
        应该将\code{std::map}视为两个模板参数\code{Key}和\code{T}上的双函子还是普罗函子（profunctor）？如何重新设计该数据类型使其满足相应特性？
\end{enumerate}